% !TEX program = xelatex
\documentclass[11pt]{article}

\usepackage[a4paper,margin=1in]{geometry}
\usepackage{amsmath,amssymb}
\usepackage{booktabs}
\usepackage{enumitem}
\usepackage{hyperref}
\usepackage{xcolor}
\usepackage{kotex}
\usepackage{listings}

\lstdefinestyle{codestyle}{
    basicstyle=\ttfamily\small,
    breaklines=true,
    frame=single,
    columns=fullflexible,
    backgroundcolor=\color{gray!5}
}
\lstset{style=codestyle}

\title{Phase 1 Code Guide\\
\large Robot Joint Motor Housing Inverse PINN}
\author{안형준}
\date{}

\begin{document}
\maketitle

\section{문서 목적}

이 문서는 \texttt{Phase1/Phase1\_Proposal\_안형준.pdf}와 \texttt{main.py}를 같이 보면서 코드 흐름을 이해하기 위한 해설 문서다.

\begin{center}
\begin{tabular}{ll}
\toprule
역할 & 파일 \\
\midrule
프로젝트 물리 모델 기준 & \texttt{Phase1/Phase1\_Proposal\_안형준.pdf} \\
현재 PINN 코드 & \texttt{main.py} \\
프로젝트 기준 메모 & \texttt{BASELINE.md} \\
\bottomrule
\end{tabular}
\end{center}

\section{프로젝트 한 줄 요약}

제안서의 물리 모델은 로봇 관절 모터 하우징을 축방향 1D 비정상 열전달 문제로 단순화한 inverse PINN이다.

\[
\text{input} = (X,\tau), \qquad
\text{output} = \theta(X,\tau)
\]

PINN은 무차원 온도장 \(\theta(X,\tau)\)를 예측하고, 동시에 inverse parameter인 \(\beta\)와 \(S\)를 학습한다. 학습 후에는 \(\beta\)와 \(S\)를 실제 물리 파라미터 \(h\), \(q_0\)로 변환한다.

\[
\beta \rightarrow h, \qquad S \rightarrow q_0
\]

\section{제안서 식과 코드의 연결}

제안서의 차원 있는 지배 방정식은 다음과 같다.

\begin{equation}
\rho c_p \frac{\partial T}{\partial t}
=
k \frac{\partial^2 T}{\partial x^2}
-
\frac{hP}{A}(T - T_\infty)
+
q_0 I(t)^2 g(x)
\end{equation}

각 항의 의미는 다음과 같다.

\begin{center}
\begin{tabular}{ll}
\toprule
항 & 의미 \\
\midrule
\(\rho c_p \frac{\partial T}{\partial t}\) & 시간에 따른 열 저장 \\
\(k \frac{\partial^2 T}{\partial x^2}\) & 축방향 전도 \\
\(-\frac{hP}{A}(T-T_\infty)\) & 외부 표면 대류 냉각 \\
\(q_0 I(t)^2 g(x)\) & 내부 Joule heating \\
\bottomrule
\end{tabular}
\end{center}

제안서에서는 학습 안정성을 위해 다음 무차원 변수를 사용한다.

\begin{equation}
X = \frac{x}{L},
\qquad
\tau = \frac{\alpha t}{L^2},
\qquad
\theta = \frac{T - T_\infty}{\Delta T_{\mathrm{ref}}},
\qquad
\alpha = \frac{k}{\rho c_p}
\end{equation}

전류도 무차원화한다.

\begin{equation}
i(\tau) = \frac{I(t)}{I_{\mathrm{ref}}}
\end{equation}

그러면 무차원 지배 방정식은 다음 형태가 된다.

\begin{equation}
\frac{\partial \theta}{\partial \tau}
=
\frac{\partial^2 \theta}{\partial X^2}
-
\beta \theta
+
S i(\tau)^2 g(X)
\end{equation}

여기서

\begin{equation}
\beta =
\frac{hPL^2}{kA}
\end{equation}

\begin{equation}
S =
\frac{q_0 I_{\mathrm{ref}}^2 L^2}
{k\Delta T_{\mathrm{ref}}}
\end{equation}

코드에서는 모든 항을 왼쪽으로 넘겨 residual 형태로 계산한다.

\begin{equation}
f_{\mathrm{PDE}}
=
\theta_\tau
-
\theta_{XX}
+
\beta \theta
-
S i(\tau)^2 g(X)
\end{equation}

PINN의 목표는 이 residual을 0에 가깝게 만드는 것이다.

\begin{equation}
f_{\mathrm{PDE}} \rightarrow 0
\end{equation}

코드 위치는 다음 함수다.

\begin{lstlisting}[language=Python]
def pde_residual(model, xtau):
\end{lstlisting}

\section{사용자가 정하는 입력값}

코드의 첫 번째 큰 구역은 사용자가 직접 정하는 값들을 모아둔 곳이다.

\begin{lstlisting}[language=Python]
# 1. User-defined inputs
\end{lstlisting}

\subsection{형상값}

\begin{center}
\begin{tabular}{lll}
\toprule
코드 변수 & 제안서 의미 & 현재 값 \\
\midrule
\texttt{L} & 모터 하우징 길이 & \(0.15~\mathrm{m}\) \\
\texttt{R\_OUTER} & 외반지름 & \(0.04~\mathrm{m}\) \\
\texttt{R\_INNER} & 내반지름 & \(0.025~\mathrm{m}\) \\
\texttt{X0} & Gaussian heat source 중심 & \(0.30\) \\
\texttt{SIGMA\_X} & Gaussian heat source 폭 & \(0.10\) \\
\texttt{SENSOR\_X} & sparse sensor 위치 & \(0.25, 0.55, 0.85\) \\
\bottomrule
\end{tabular}
\end{center}

제안서에서 열원 중심과 폭은 다음과 같이 설정했다.

\[
x_0 = 0.3L,
\qquad
\sigma = 0.1L
\]

무차원 좌표 \(X=x/L\)를 쓰면 코드에서는 다음처럼 쓴다.

\[
X_0 = 0.30,
\qquad
\sigma_X = 0.10
\]

\subsection{재료 물성치}

\begin{center}
\begin{tabular}{lll}
\toprule
코드 변수 & 제안서 의미 & 현재 값 \\
\midrule
\texttt{RHO} & 밀도 \(\rho\) & \(2700~\mathrm{kg/m^3}\) \\
\texttt{CP} & 비열 \(c_p\) & \(900~\mathrm{J/(kg\cdot K)}\) \\
\texttt{K} & 열전도율 \(k\) & \(167~\mathrm{W/(m\cdot K)}\) \\
\texttt{T\_INF} & 주변 공기 온도 \(T_\infty\) & \(25^\circ\mathrm{C}\) \\
\bottomrule
\end{tabular}
\end{center}

\subsection{무차원화 기준값}

\begin{center}
\begin{tabular}{lll}
\toprule
코드 변수 & 의미 & 현재 값 \\
\midrule
\texttt{DELTA\_T\_REF} & 온도 기준값 & \(50^\circ\mathrm{C}\) \\
\texttt{I\_REF} & 전류 기준값 & \(5.0~\mathrm{A}\) \\
\texttt{T\_FINAL} & 최종 해석 시간 & \(1200~\mathrm{s}\) \\
\bottomrule
\end{tabular}
\end{center}

\subsection{Synthetic data용 true parameter}

\begin{center}
\begin{tabular}{lll}
\toprule
코드 변수 & 의미 & 현재 값 \\
\midrule
\texttt{H\_TRUE} & synthetic data 생성용 실제 \(h\) & \(40~\mathrm{W/(m^2\cdot K)}\) \\
\texttt{Q0\_TRUE} & synthetic data 생성용 실제 \(q_0\) & \(8.0\times10^4~\mathrm{W/(m^3 A^2)}\) \\
\bottomrule
\end{tabular}
\end{center}

이 두 값은 inverse PINN이 나중에 맞혀야 할 정답 역할을 한다.

\section{자동 계산값}

입력값에서 자동 계산되는 값들은 다음 구역에 있다.

\begin{lstlisting}[language=Python]
# 2. Derived values
\end{lstlisting}

\begin{center}
\begin{tabular}{lll}
\toprule
코드 변수 & 공식 & 의미 \\
\midrule
\texttt{ALPHA} & \(\alpha = k/(\rho c_p)\) & 열확산계수 \\
\texttt{AREA} & \(A=\pi(R_o^2-R_i^2)\) & 축방향 전도 단면적 \\
\texttt{PERIMETER} & \(P=2\pi R_o\) & 외부 대류 둘레 \\
\texttt{TAU\_FINAL} & \(\tau_f=\alpha t_f/L^2\) & 무차원 최종 시간 \\
\texttt{BETA\_TRUE} & \(\beta_{\mathrm{true}}=h_{\mathrm{true}}PL^2/(kA)\) & synthetic data용 true \(\beta\) \\
\texttt{S\_TRUE} & \(S_{\mathrm{true}}=q_{0,\mathrm{true}}I_{\mathrm{ref}}^2L^2/(k\Delta T_{\mathrm{ref}})\) & synthetic data용 true \(S\) \\
\bottomrule
\end{tabular}
\end{center}

여기 값들은 보통 직접 수정하지 않고, 앞의 입력값을 바꿔서 간접적으로 바꾼다.

\section{차원/무차원 변환 함수}

\subsection{시간 변환}

제안서의 시간 무차원화는 다음과 같다.

\begin{equation}
\tau = \frac{\alpha t}{L^2}
\end{equation}

코드 함수는 다음 두 개다.

\begin{lstlisting}[language=Python]
tau_from_time(t_seconds)
time_from_tau(tau)
\end{lstlisting}

\subsection{\texorpdfstring{\(h\)와 \(\beta\) 변환}{h and beta conversion}}

\begin{equation}
\beta = \frac{hPL^2}{kA}
\end{equation}

\begin{equation}
h = \frac{\beta kA}{PL^2}
\end{equation}

코드 함수

\begin{lstlisting}[language=Python]
beta_from_physical_h(h_value)
physical_h_from_beta(beta)
\end{lstlisting}

\subsection{\texorpdfstring{\(q_0\)와 \(S\) 변환}{q0 and S conversion}}

\begin{equation}
S =
\frac{q_0 I_{\mathrm{ref}}^2 L^2}
{k\Delta T_{\mathrm{ref}}}
\end{equation}

\begin{equation}
q_0 =
\frac{S k \Delta T_{\mathrm{ref}}}
{I_{\mathrm{ref}}^2 L^2}
\end{equation}

코드 함수

\begin{lstlisting}[language=Python]
s_from_physical_q0(q0_value)
physical_q0_from_s(s_value)
\end{lstlisting}

\section{Gaussian heat source}

제안서의 열원 분포는 Gaussian이다.

\begin{equation}
g(X)
=
\exp\left[
-
\frac{(X-X_0)^2}{2\sigma_X^2}
\right]
\end{equation}

코드 위치

\begin{lstlisting}[language=Python]
def gaussian_heat_source(x):
\end{lstlisting}

\begin{center}
\begin{tabular}{ll}
\toprule
위치 & 의미 \\
\midrule
\(X=0.30\) 근처 & winding 근처, 열 발생 큼 \\
\(X=0.85\) 근처 & heat source에서 멀어서 열 발생 작음 \\
\bottomrule
\end{tabular}
\end{center}

\section{Current profile}

코드 위치

\begin{lstlisting}[language=Python]
def current_profile(tau):
\end{lstlisting}

현재 사용자가 지정한 전류 입력은 두 번 켜고 끄는 duty cycle이다.

\begin{center}
\begin{tabular}{lll}
\toprule
구간 & 상태 & 의미 \\
\midrule
\(0 \le t < 300~\mathrm{s}\) & ON & Joule heating 작동 \\
\(300 \le t < 600~\mathrm{s}\) & OFF & 냉각 중심 \\
\(600 \le t < 900~\mathrm{s}\) & ON & Joule heating 재작동 \\
\(900 \le t \le 1200~\mathrm{s}\) & OFF & 냉각 중심 \\
\bottomrule
\end{tabular}
\end{center}

inverse problem에서 이 구조가 중요한 이유는 다음과 같다.

\begin{center}
\begin{tabular}{lll}
\toprule
구간 & 물리 의미 & 추정에 도움 되는 값 \\
\midrule
current ON & Joule heating 항 존재 & \(q_0\), \(S\) \\
current OFF & Joule heating 항 사라짐 & \(h\), \(\beta\) \\
\bottomrule
\end{tabular}
\end{center}

\section{Synthetic sensor data}

코드 위치

\begin{lstlisting}[language=Python]
def synthetic_theta_no_conduction(x, tau):
def make_synthetic_sensor_data(device):
\end{lstlisting}

현재 synthetic data는 먼저 단순한 local ODE로 생성한다.

\begin{equation}
\frac{d\theta}{d\tau}
=
-\beta_{\mathrm{true}}\theta
+
S_{\mathrm{true}}i(\tau)^2g(X)
\end{equation}

즉, 현재 synthetic data에는 축방향 전도항 \(\theta_{XX}\)가 아직 없다.

\begin{equation}
\theta_{XX} \quad \text{not included yet}
\end{equation}

현재 이렇게 둔 이유는 다음과 같다.

\begin{enumerate}[leftmargin=2em]
\item inverse PINN 구조가 돌아가는지 먼저 확인
\item sensor data loss 연결 확인
\item 나중에 full PDE finite difference solver로 교체
\end{enumerate}

현재 sensor data 개수는 다음과 같다.

\[
3~\text{sensors} \times 41~\text{time points}
=
123~\text{points}
\]

\section{신경망 구조}

코드 위치

\begin{lstlisting}[language=Python]
class MotorHousingPINN(nn.Module):
\end{lstlisting}

신경망 입력과 출력은 다음과 같다.

\[
(X,\tau) \longrightarrow \theta(X,\tau)
\]

현재 네트워크 구조

\begin{lstlisting}
Linear(2 -> 32)
Tanh
Linear(32 -> 32)
Tanh
Linear(32 -> 32)
Tanh
Linear(32 -> 1)
\end{lstlisting}

모델 내부에는 trainable parameter 두 개가 들어 있다.

\begin{lstlisting}[language=Python]
self.raw_beta
self.raw_s
\end{lstlisting}

코드는 바로 \(\beta\), \(S\)를 학습하지 않고 raw parameter를 학습한 뒤 softplus를 통과시킨다.

\begin{equation}
\beta = \mathrm{softplus}(\texttt{raw\_beta})
\end{equation}

\begin{equation}
S = \mathrm{softplus}(\texttt{raw\_s})
\end{equation}

이렇게 하면 \(\beta\)와 \(S\)가 항상 양수가 된다. 물리적으로 \(h\), \(q_0\), \(\beta\), \(S\)가 음수면 이상하기 때문에 이 처리를 둔다.

\section{점 샘플링}

PINN은 전체 격자를 고정해서 쓰기보다, 여러 종류의 점을 뽑아서 loss를 계산한다.

\begin{center}
\begin{tabular}{lll}
\toprule
함수 & 뽑는 점 & 제안서 loss와 연결 \\
\midrule
\texttt{sample\_collocation\_points()} & PDE 검사점 & \(L_{\mathrm{PDE}}\) \\
\texttt{sample\_initial\_points()} & \(\tau=0\) 점 & \(L_{\mathrm{IC}}\) \\
\texttt{sample\_boundary\_points()} & \(X=0, X=1\) 점 & \(L_{\mathrm{BC}}\) \\
\texttt{make\_synthetic\_sensor\_data()} & 센서 위치와 시간점 & \(L_{\mathrm{data}}\) \\
\bottomrule
\end{tabular}
\end{center}

\section{자동미분}

코드 위치

\begin{lstlisting}[language=Python]
def gradient(output, inputs):
\end{lstlisting}

PINN에서는 신경망 출력 \(\theta\)를 직접 미분해서 PDE residual을 만든다.

필요한 미분값은 다음과 같다.

\begin{equation}
\theta_\tau =
\frac{\partial \theta}{\partial \tau}
\end{equation}

\begin{equation}
\theta_X =
\frac{\partial \theta}{\partial X}
\end{equation}

\begin{equation}
\theta_{XX} =
\frac{\partial^2 \theta}{\partial X^2}
\end{equation}

코드에서 이 값들을 계산하는 부분

\begin{lstlisting}[language=Python]
grad_theta = gradient(theta, xtau)
theta_x = grad_theta[:, 0:1]
theta_tau = grad_theta[:, 1:2]

grad_theta_x = gradient(theta_x, xtau)
theta_xx = grad_theta_x[:, 0:1]
\end{lstlisting}

\section{Loss 함수}

\subsection{PDE loss}

코드 위치

\begin{lstlisting}[language=Python]
pde_residual()
pde_loss()
\end{lstlisting}

제안서의 PDE loss는 다음과 같은 형태다.

\begin{equation}
L_{\mathrm{PDE}}
=
\frac{1}{N_f}
\sum_{i=1}^{N_f}
f_{\mathrm{PDE}}(X_i,\tau_i)^2
\end{equation}

\subsection{Initial condition loss}

제안서의 초기조건

\begin{equation}
T(x,0)=T_\infty
\end{equation}

무차원 온도에서는 다음과 같다.

\begin{equation}
\theta(X,0)=0
\end{equation}

코드 위치

\begin{lstlisting}[language=Python]
initial_condition_loss()
\end{lstlisting}

\subsection{Boundary condition loss}

제안서의 단열 경계조건

\begin{equation}
\frac{\partial T}{\partial x}(0,t)=0
\end{equation}

\begin{equation}
\frac{\partial T}{\partial x}(L,t)=0
\end{equation}

무차원 좌표에서는 다음과 같다.

\begin{equation}
\frac{\partial \theta}{\partial X}(0,\tau)=0
\end{equation}

\begin{equation}
\frac{\partial \theta}{\partial X}(1,\tau)=0
\end{equation}

코드 위치

\begin{lstlisting}[language=Python]
boundary_condition_loss()
\end{lstlisting}

\subsection{Sensor data loss}

제안서의 sparse sensor data loss는 다음 형태다.

\begin{equation}
L_{\mathrm{data}}
=
\frac{1}{N_s}
\sum_{i=1}^{N_s}
\left[
T_\phi(x_i,t_i)-T_{\mathrm{sensor},i}
\right]^2
\end{equation}

현재 코드는 무차원 온도 \(\theta\)를 쓰므로 다음 형태로 계산한다.

\begin{equation}
L_{\mathrm{data}}
=
\frac{1}{N_s}
\sum_{i=1}^{N_s}
\left[
\theta_\phi(X_i,\tau_i)-\theta_{\mathrm{sensor},i}
\right]^2
\end{equation}

코드 위치

\begin{lstlisting}[language=Python]
sensor_data_loss()
\end{lstlisting}

\section{main 함수}

코드 위치

\begin{lstlisting}[language=Python]
def main():
\end{lstlisting

현재 \texttt{main()}은 학습을 오래 돌리는 함수가 아니라 smoke test용이다.

현재 역할

\begin{enumerate}[leftmargin=2em]
\item device 선택
\item 모델 생성
\item PDE loss 계산
\item IC loss 계산
\item BC loss 계산
\item synthetic sensor data loss 계산
\item \(\beta\), \(S\)를 \(h\), \(q_0\)로 변환해서 출력
\item 코드 구조가 깨지지 않는지 확인
\end{enumerate}

출력값 의미

\begin{center}
\begin{tabular}{ll}
\toprule
출력 & 의미 \\
\midrule
\texttt{tau\_final} & \(1200~\mathrm{s}\)를 무차원 시간으로 바꾼 값 \\
\texttt{beta\_true} & \(h_{\mathrm{true}}\)에서 계산한 정답 \(\beta\) \\
\texttt{S\_true} & \(q_{0,\mathrm{true}}\)에서 계산한 정답 \(S\) \\
\texttt{beta guess} & 학습 전 모델의 \(\beta\) 초기값 \\
\texttt{S guess} & 학습 전 모델의 \(S\) 초기값 \\
\texttt{PDE loss} & PDE residual 오차 \\
\texttt{IC loss} & 초기조건 오차 \\
\texttt{BC loss} & 단열 경계조건 오차 \\
\texttt{Data loss} & sensor data 오차 \\
\bottomrule
\end{tabular}
\end{center}

\section{현재 코드의 한계}

현재 코드가 아직 하지 않는 것

\begin{enumerate}[leftmargin=2em]
\item optimizer를 이용한 실제 학습
\item full PDE synthetic data 생성
\item finite difference reference solver
\item \(\beta\), \(S\)가 true value로 수렴하는지 확인
\item temperature contour plot
\item sensor data fitting plot
\item parameter convergence plot
\end{enumerate}

가장 중요한 한계는 다음 함수다.

\begin{lstlisting}[language=Python]
synthetic_theta_no_conduction()
\end{lstlisting}

이 함수는 아직 전도항을 포함하지 않는다.

\[
\theta_{XX} \quad \text{is not included yet}
\]

따라서 최종 reference data generator가 아니라 첫 연결 확인용 synthetic data generator로 봐야 한다.

\section{다음 구현 순서}

추천 순서

\begin{enumerate}[leftmargin=2em]
\item training loop 추가
\item total loss weight 정리
\item \(\beta\), \(S\) 학습이 되는지 짧게 확인
\item synthetic data generator를 full PDE finite difference solver로 교체
\item 결과 plot 추가
\item learned \(\beta\), \(S\)를 \(h\), \(q_0\)로 변환해서 보고
\end{enumerate}

가장 먼저 볼 코드 지점

\begin{lstlisting}[language=Python]
# 1. User-defined inputs
\end{lstlisting}

그 다음 볼 코드 지점

\begin{lstlisting}[language=Python]
def pde_residual(model, xtau):
\end{lstlisting}

제안서의 핵심 PDE가 실제로 코드로 들어간 위치이기 때문이다.

\end{document}
